\section*{Erd\H{o}s Problem 1144 --- GPT 6 Astra Ultra partial attempt}


\subsection*{FORMAL RESTATEMENT}
Independently assign each prime $p$ a sign $f(p)\in\{-1,1\}$ with equal probabilities, and extend $f$ completely multiplicatively to positive integers. Does
\[
\limsup_{N\to\infty}\frac{1}{\sqrt N}\sum_{m\leq N}f(m)=+\infty
\]
hold with probability one?

\subsection*{CONTEXT AND SCOPE}
This replaces the earlier statement-only record with an elementary mathematical attempt. The arguments below establish only their stated partial results; no novelty or full solution is claimed.

Complete multiplicativity is essential: this is not the squarefree-supported random multiplicative model.

\subsection*{ATTACK PLAN AND WORK}
Write $S_N=\sum_{n\le N}f(n)$. Independence of the prime signs implies
$\mathbb E f(n)=1$ for a square and zero otherwise, and
$\mathbb E[f(m)f(n)]=1$ precisely when $mn$ is a square. Grouping integers
by their squarefree parts gives the exact identities
\[
\mathbb E S_N=\lfloor\sqrt N\rfloor,\qquad
\mathbb E S_N^2=\sum_{\substack{a\le N\\a\ {
m squarefree}}}
                  \lfloor\sqrt{N/a}\rfloor^2.\tag{1}
\]
Here a squarefree part $a$ contributes pairs $(au^2,av^2)$.

These identities already force logarithmically growing second moments after
normalization. Let $Q(t)$ count squarefree positive integers at most $t$.
The union bound over all possible square divisors gives, for integral $t$,
$Q(t)\ge t-t\sum_{j\ge2}j^{-2}\ge t/4$, since
$\sum_{j\ge2}j^{-2}\le1/4+\int_2^\infty x^{-2}\,dx=3/4$.
Partial summation consequently gives
$\sum_{a\le N,\ a\ {
m squarefree}}1/a\ge H_N/4$.
Using $\lfloor y\rfloor\ge y/2$ for $y\ge1$ in (1),
\[NH_N/16\le\mathbb E S_N^2\le NH_N.\tag{2}\]
In particular the variance is unbounded relative to $N$.

This is not an almost-sure positive-limsup argument. Large second moments can
come from rare events, and a signed random variable need not have a large
positive tail merely because its variance is large. Temporal dependence also
prevents applying independent-event Borel--Cantelli to a sequence of $S_N$.

\subsection*{FINITE COMPUTATIONAL CHECK}
Every assignment of prime signs was enumerated for each $1\le N\le16$. Exact rational averages agree with both moment identities and satisfy the harmonic-number bounds.

\subsection*{VERIFICATION AND REMAINING GAP}
Obtain quantitative positive-tail events at infinitely many scales with
enough control of their dependence. Moment growth alone does not establish
the almost-sure one-sided assertion.

\subsection*{SOURCES}
https://www.erdosproblems.com/1144

https://www.erdosproblems.com/history/1144

\subsection*{FINAL}
\textbf{UNRESOLVED.} The full source problem remains outside the scope of the proved partial results.

\subsection*{COMPLETION ESTIMATE}
$8\%$. This is a conservative subjective estimate toward the whole problem, not a probability that the argument is correct or a claim that a fixed fraction of the problem has been solved.
